\chapter{Arithmetic Virasoro Positivity Bootstrap}
\label{v4-arith:virasoro-bootstrap}

\emph{With the determinant form of
Chapter~\ref{v4-arith:rh-reformulation} in place, the
Virasoro-polynomial sequence
\(T_{n}^{\mathrm{Vir}}=(L_{0}^{\mathrm{Vir}})^{n}-(L_{0}^{\mathrm{Vir}}-\mathrm{id})^{n}\)
is not the Li operator. It proves finite-sector Fock positivity, but
it does not construct a Hilbert--P\'olya operator and it does not give
RH. The RH-strength input is the positive Hilbert--P\'olya trace of
\(\mathcal L_{n}=\mathrm{id}-(\mathrm{id}-\Theta_{\xi}^{-1})^{n}\)
from Chapter~\ref{v4-arith:rh-reformulation}.}

\section{Statement Recalled}
\label{v4-arith:vb:statement}

\subsection{The Operator Sequence}
\label{v4-arith:vb:operator-sequence}

\begin{definition}[formal Virasoro polynomial sequence]
\label{v4-arith:vb:def:T_n}
\ClaimStatusDefinitional. Let \(\mathcal{V}^{\mathrm{prim}}_{\mathrm{Heis},\mathrm{prim}}\)
denote the primary (non-descendant) subspace of the arithmetic
Heisenberg Fock, equipped with the Petersson inner product
\(\langle\,\cdot\,,\,\cdot\,\rangle^{\mathrm{Pet}}\) of
\Cref{v4-arith:thm:P3-resolution}. For every integer \(n\geq 1\)
define
\begin{equation}
\label{v4-arith:vb:eq:T_n-def}
T_{n}^{\mathrm{Vir}}\;:=\;(L_{0}^{\mathrm{Vir}})^{n}
\,-\,(L_{0}^{\mathrm{Vir}}-\mathrm{id})^{n}
\;\in\;\mathrm{End}\bigl(\mathcal{V}^{\mathrm{prim}}_{\mathrm{Heis},\mathrm{prim}}\bigr).
\end{equation}
Here \(L_{0}^{\mathrm{Vir}}\) is the Fock/Virasoro conformal-energy
operator. It is not the Hilbert--P\'olya determinant-trace operator
\(H_{\mathrm{HP}}\), the Bost--Connes Hamiltonian \(N_{\mathrm{BC}}\),
or the zero-divisor operator \(\Theta_{\xi}=1/2+iH_{\mathrm{HP}}\). Thus
\(T_{n}^{\mathrm{Vir}}\) is a legitimate Fock polynomial on any
common Virasoro spectral core, but it is not Li's transform.
In the sequel, later occurrences of
\(T_n\) mean \(T_{n}^{\mathrm{Vir}}\) unless the Hilbert--P\'olya
Li transform \(\mathcal L_n\) is written explicitly.
\end{definition}

\begin{remark}[why \(T_{n}^{\mathrm{Vir}}\) is not the Li object]
\label{v4-arith:vb:rem:why-T_n}
Two conventions for Li's criterion circulate in the literature. The
original Li 1997 J.~Number Theory paper defines
\(\lambda_{n}=\sum_{\rho}(1-(1-1/\rho)^{n})\); Bombieri--Lagarias J.
Number Theory 77 (1999) uses the equivalent closed form
\(\lambda_{n}=(n-1)!^{-1}(d/ds)^{n}[s^{n-1}\log\xi(s)]\big|_{s=1}\).
The corresponding operator is
\(\mathcal L_{n}=\mathrm{id}-(\mathrm{id}-\Theta_{\xi}^{-1})^{n}\),
with \(\Theta_{\xi}\) the normal zero-divisor operator of
\Cref{v4-arith:rh:hyp:HPAr}. A polynomial
\((L_{0}^{\mathrm{Vir}})^{n}-(L_{0}^{\mathrm{Vir}}-\mathrm{id})^{n}\)
does not become Li's transform by regularisation. Consequently, every
claim below involving \(T_{n}^{\mathrm{Vir}}\) is a formal
Fock/Virasoro claim unless it is explicitly restated in terms of
\(\mathcal L_{n}\) and the positive Weil--Connes trace.
\end{remark}

\subsection{Conditional Chain: VB\(^{*}\) \(\Rightarrow\) Li Positivity \(\Leftrightarrow\) RH}
\label{v4-arith:vb:equivalence-chain}

\begin{proposition}[formal Virasoro positivity is not Li positivity]
\label{v4-arith:vb:thm:three-equivalences}
\ClaimStatusProvedHere. Positivity of the formal Virasoro polynomial
\(T_{n}^{\mathrm{Vir}}\) on finite Fock sectors does not imply
\(\lambda_{n}\geq0\), and therefore does not imply RH. The
one-directional RH implication in Vol~IV is instead
\[
\mathrm{HP\text{--}Li}\Longrightarrow \lambda_{n}\geq 0\ \forall n
\Longleftrightarrow \mathrm{RH},
\]
where HP--Li is the positive trace statement of
\Cref{v4-arith:rh:thm:VB-implies-F-RH}.
\end{proposition}

\begin{proof}
Li's criterion is Li 1997 J.~Number Theory, reproduced in
Bombieri--Lagarias J. Number Theory 77 (1999). The Li summand is
\(1-(1-1/\rho)^{n}\). The spectral transform realising this summand
is therefore \(z\mapsto1-(1-z^{-1})^{n}\) applied to the zero
operator \(z=\Theta_{\xi}\). The polynomial \(x^{n}-(x-1)^{n}\)
applied to a Virasoro energy \(x=L_{0}^{\mathrm{Vir}}\) is a different
function of a different operator. Positivity of the latter is a Fock
calculation; it has no trace identity with Li's coefficients unless
the Hilbert--P\'olya trace identity has already been replaced by the
correct \(\Theta_{\xi}\)-transform. The HP--Li implication follows
from \Cref{v4-arith:rh:thm:VB-implies-F-RH}.
\end{proof}

\begin{remark}[what is genuinely equivalent]
\label{v4-arith:vb:rem:equiv-domain}
The known RH-equivalent positivity languages are Li positivity, Weil
positivity for the number-field explicit formula, and the positive
zero-measure conjecture for \(\Theta_{\xi}\),
VB\(^{*}_{\mathrm{core}}\). Under explicit hypotheses these languages
give the same positivity condition. None of the RH-equivalent
positivity statements is unconditional here.
\end{remark}

\begin{remark}[the level at which RH lives]
\label{v4-arith:vb:rem:rh-level}
RH lives on the zero-divisor operator
\(\Theta_{\xi}=1/2+iH_{\mathrm{HP}}\), not on the Virasoro energy
operator. The Bost--Connes Hamiltonian and the Virasoro \(L_{0}\)
give Euler weights and Fock energies; neither carries
\(\mathrm{div}\,\xi_{\mathbb R}\). A real spectral measure appears only after
H-P\(^{\mathrm{Ar}}\) supplies \(H_{\mathrm{HP}}\) and identifies
\(\det_{\infty}(s-\Theta_{\xi})\) with \(\xi_{\mathbb R}(s)\).
VB\(^{*}\), in its RH-strength form, is therefore a positivity
statement for the Weil--Connes trace of functions of
\(\Theta_{\xi}\), together with determinacy of the corresponding
zero measure.
\end{remark}

\section{Obstructions}
\label{v4-arith:vb:table}

\subsection{Virasoro levels}
\label{v4-arith:vb:cycle:A1}

For \(n=1\) one has \(T_{1}^{\mathrm{Vir}}=\mathrm{id}\). For
\(n=2\) one has
\[
(L_{0}^{\mathrm{Vir}})^{2}-(L_{0}^{\mathrm{Vir}}-\mathrm{id})^{2}
=2L_{0}^{\mathrm{Vir}}-\mathrm{id}.
\]
Thus \(T_{2}^{\mathrm{Vir}}\succeq0\) is the Fock threshold
\(L_{0}^{\mathrm{Vir}}\succeq\frac12\mathrm{id}\). It is not a
statement about zeta zeros. Li's coefficient at level \(2\) is the
trace of \(z\mapsto1-(1-z^{-1})^{2}\) on \(\Theta_{\xi}\), and equality
between these two functional calculi would itself be an additional
Hilbert--P\'olya theorem.

\subsection{Moment problem}
\label{v4-arith:vb:cycle:A2}

The Hamburger moment problem belongs to the regularised spectral
measure of \(\Theta_{\xi}\). There is no identity
\[
\mathrm{Tr}_{\mathrm{reg}}(T_{n}^{\mathrm{Vir}})=n\lambda_{n}.
\]
The valid moment sequence is obtained from the zero-divisor operator and the
Weil--Connes trace. Stieltjes' theorem supplies the algebraic test:
a sequence is the moment sequence of a positive measure on
\([0,\infty)\) precisely when all Hankel and shifted Hankel minors are
non-negative.

\subsection{Kac determinant}
\label{v4-arith:vb:cycle:A3}

At \(c=2\), the Kac parametrisation gives
\[
1-\frac{6}{m(m+1)}=2,\qquad m(m+1)=-6.
\]
The null weights form the countable set
\(\{h_{r,s}(2):r,s\geq1\}\). Outside this set the contravariant form
is non-degenerate. Non-degeneracy is weaker than positivity; the
Petersson sign is not determined by the non-vanishing of the Kac
determinant.

\subsection{Fock combinatorics}
\label{v4-arith:vb:cycle:A4}

On the free-boson number state
\(|\mathbf n\rangle\), \(L_{0}^{\mathrm{Vir}}\) has eigenvalue
\[
m(\mathbf n)=\sum_{i\geq1}i n_i .
\]
For \(m(\mathbf n)\geq1\),
\[
T_{n}^{\mathrm{Vir}}|\mathbf n\rangle
=\bigl(m(\mathbf n)^n-(m(\mathbf n)-1)^n\bigr)|\mathbf n\rangle
\]
has non-negative eigenvalue. The vacuum has \(m=0\); for even \(n\)
the eigenvalue is negative. Hence the Fock positivity statement is a
non-vacuum statement.

\subsection{Weil trace}
\label{v4-arith:vb:cycle:A5}

Weil positivity is a positivity theorem for the explicit formula and
is equivalent to RH over number fields. In the Vol~IV notation its
operator representative is the positive regularised trace of
\(\mathcal L_n(\Theta_{\xi})\), not a trace of
\(T_n^{\mathrm{Vir}}\). Koszul--Petersson compatibility can identify
the Weil trace pairing with a Petersson pairing on admissible
arithmetic states; it cannot identify the Virasoro energy spectrum
with the zero divisor.

\subsection{de Branges datum}
\label{v4-arith:vb:cycle:A6}

For a Hermite--Biehler function \(E\), the de Branges kernel is
\[
K(w,z)=\frac{E(w)\overline{E(z)}-E^{*}(w)\overline{E^{*}(z)}}
{2\pi i(\overline z-w)}.
\]
The real entire function
\(\Xi(z)=\xi_{\mathbb R}(1/2+iz)\) satisfies \(\Xi^{*}=\Xi\), so
\(\Xi\) itself is not a Hermite--Biehler function. A de Branges
reformulation of RH requires a genuine datum \(E\) whose divisor
controls the divisor of \(\Xi\). The Hermite--Biehler condition is then
part of the RH-strength input, not a consequence of formal Virasoro
positivity.

\section{Level-by-Level Positivity}
\label{v4-arith:vb:level-closure}

\subsection{Level \(n=1\): Unconditional}
\label{v4-arith:vb:n=1}

\begin{theorem}[\(n=1\) arithmetic Virasoro positivity]
\label{v4-arith:vb:thm:n=1}
\ClaimStatusProvedHere. \(T_{1}^{\mathrm{Vir}}=(L_{0}^{\mathrm{Vir}})-(L_{0}^{\mathrm{Vir}}-\mathrm{id})=\mathrm{id}\).
Hence \(T_{1}\succeq 0\) unconditionally on
\(\mathcal{V}^{\mathrm{prim}}_{\mathrm{Heis},\mathrm{prim}}\).
\end{theorem}

\begin{proof}
Direct computation. The operator is the identity.
\end{proof}

\begin{remark}[numerical consistency]
\label{v4-arith:vb:rem:n=1-consistency}
The ordinary trace of the identity on
\(\mathcal{V}^{\mathrm{prim}}_{\mathrm{Heis},\mathrm{prim}}\) is not
defined without a trace-class or regularized-trace convention.
Bombieri--Lagarias \S2 gives
\(\lambda_{1}=\tfrac{1}{2}(2+\gamma-\log(4\pi))\approx 0.02310\). The
positivity of this number is consistent but not itself a proof of
\(T_{1}\succeq 0\); the proof is the identity-operator identity.
The number \(\lambda_{1}\) has been known to be positive since Li 1997.
\end{remark}

\subsection{Level \(n=2\): Formal Threshold}
\label{v4-arith:vb:n=2}

\begin{theorem}[\(n=2\) Virasoro threshold]
\label{v4-arith:vb:thm:n=2}
\ClaimStatusProvedHere. \(T_{2}^{\mathrm{Vir}}=2L_{0}^{\mathrm{Vir}}-\mathrm{id}\).
Hence \(T_{2}^{\mathrm{Vir}}\succeq 0\) on
\(\mathcal{V}^{\mathrm{prim}}_{\mathrm{Heis},\mathrm{prim}}\) if and
only if \(L_{0}^{\mathrm{Vir}}\succeq\tfrac{1}{2}\mathrm{id}\) on the
same space.
\end{theorem}

\begin{proof}
Direct expansion:
\(T_{2}^{\mathrm{Vir}}=(L_{0}^{\mathrm{Vir}})^{2}-(L_{0}^{\mathrm{Vir}}-\mathrm{id})^{2}
=2L_{0}^{\mathrm{Vir}}-\mathrm{id}\). The inequality
\(T_{2}^{\mathrm{Vir}}\succeq 0\) is therefore equivalent to
\(2L_{0}^{\mathrm{Vir}}\succeq\mathrm{id}\), i.e.~to
\(L_{0}^{\mathrm{Vir}}\succeq\tfrac{1}{2}\mathrm{id}\).
\end{proof}

\begin{corollary}[\(n=2\) is not a zero-placement theorem]
\label{v4-arith:vb:cor:n=2-is-RH}
\ClaimStatusProvedHere. The implication
\[
T_{2}^{\mathrm{Vir}}\succeq0\quad\Longrightarrow\quad
\mathrm{Re}(\rho)=1/2\text{ for every non-trivial zero }\rho
\]
is false without an additional Hilbert--P\'olya identification of the
Virasoro energy with the zero-divisor operator. H-P\(^{\mathrm{Ar}}\) supplies
\(\Theta_{\xi}\), not an equality
\(L_{0}^{\mathrm{Vir}}=\Theta_{\xi}\).
\end{corollary}

\begin{proof}
The spectrum of \(L_{0}^{\mathrm{Vir}}\) is the conformal-energy
spectral measure of the Fock representation. The operator
\(\Theta_{\xi}\) is constrained by the determinant divisor
\(\mathrm{div}\,\xi_{\mathbb R}\). Equality of these data would
identify a non-negative Fock energy measure with a complex divisor. No
such equality follows from P1, P2, P3, or the definition of
\(T_{2}^{\mathrm{Vir}}\). The valid RH implication uses
the Li transform \(\mathcal L_{n}\) of \(\Theta_{\xi}\), as in
\Cref{v4-arith:rh:prop:Li-is-Virasoro}.
\end{proof}

\begin{remark}[the spectral barrier]
\label{v4-arith:vb:rem:n=2-is-barrier}
The \(n=2\) Virasoro calculation is a threshold calculation inside
the Fock representation. The RH barrier is the construction of the
zero-divisor operator \(\Theta_{\xi}\) and of a positive trace on its Li
transforms. Confusing these two operators turns Fock positivity into
a false Hilbert--P\'olya argument.
\end{remark}

\subsection{Levels \(n=3,4\): Bombieri-Lagarias Upper Bound}
\label{v4-arith:vb:n=3-4}

\begin{theorem}[Bombieri--Lagarias asymptotic bound]
\label{v4-arith:vb:thm:BL-bound}
\ClaimStatusProvedElsewhere (Bombieri--Lagarias J. Number Theory
77 (1999):274--287, Theorem 2). For every \(n\geq 1\),
\begin{equation}
\label{v4-arith:vb:eq:BL-bound}
\lambda_{n}\;\leq\;\frac{n(c_{0}+\log n)}{2},
\qquad c_{0}:=2+\gamma-\log(4\pi)\approx 0.04619.
\end{equation}
The matching lower bound \(\lambda_{n}\geq -\tfrac{1}{2}n\log n-O(n)\)
holds unconditionally.
\end{theorem}

\begin{corollary}[finite Li data are scalar data]
\label{v4-arith:vb:cor:n=3-4-trace}
\ClaimStatusProvedHere. For any finite set \(F\subset\mathbb N\), the
inequalities \(\lambda_{n}\geq0\) for \(n\in F\) are scalar
inequalities for the Li transforms
\(\mathcal L_{n}=1-(1-\Theta_{\xi}^{-1})^{n}\). They do not imply
operator positivity of \(T_{n}^{\mathrm{Vir}}\), and they do not
imply RH. Bombieri--Lagarias bounds and finite numerical evaluations
are therefore finite scalar controls, not Virasoro positivity
theorems.
\end{corollary}

\begin{proof}
Li's criterion requires \(\lambda_{n}\geq0\) for every \(n\geq1\).
A finite set of inequalities is therefore not equivalent to RH. The
identity \(\lambda_{n}=\mathrm{Tr}_{\mathrm{reg}}(\mathcal L_{n})\) is
a scalar trace identity for \(\Theta_{\xi}\). A positive trace of one operator
does not force positivity of another operator: a diagonal matrix with
eigenvalues \(-1\) and \(2\) has positive trace and is not positive.
The Virasoro polynomial \(T_{n}^{\mathrm{Vir}}\) is a different
functional calculus from \(\mathcal L_{n}\).
\end{proof}

\begin{remark}[the trace-positivity gap]
\label{v4-arith:vb:rem:trace-gap}
The number-theoretic trace is attached to
\(\mathcal L_{n}(\Theta_{\xi})\), not to \(T_{n}^{\mathrm{Vir}}\).
Li positivity for all \(n\) is equivalent to RH. Operator positivity
of a Virasoro polynomial is a separate Fock statement, and a scalar
trace cannot bridge the two without a positive, determinate
Weil--Connes representation of the zero measure.
\end{remark}

\section{Moment Problem Reformulation}
\label{v4-arith:vb:moment}

\subsection{The Shifted Zero Measure}
\label{v4-arith:vb:shifted-measure}

\begin{definition}[off-line defect measure]
\label{v4-arith:vb:def:zero-measure}
\ClaimStatusDefinitional. Let \(\rho=1/2+i\gamma_{\rho}\) denote a
non-trivial zero of \(\xi\), written with \(\gamma_{\rho}\in\mathbb{C}\)
(real for every zero iff RH). Under the symmetrisation
\(\rho\leftrightarrow 1-\overline{\rho}\), the pair set
\(Z=\{\rho,1-\overline{\rho}\}\) pushes to the off-line defect measure
\(\nu_{\mathrm{def}}\) on \(\mathbb{R}_{\geq0}\) via
\begin{equation}
\int f\,d\nu_{\mathrm{def}}\;:=\;\sum_{\rho}f(|\mathrm{Im}\,\gamma_{\rho}|^{2})
\end{equation}
for test functions \(f\) of suitable growth, Bost--Connes-regularised.
The measure \(\nu_{\mathrm{def}}\) is supported at \(0\) exactly when
RH holds. It is a defect detector, not Li's positive moment measure.
\end{definition}

\begin{remark}
Under RH, \(\gamma_{\rho}\in\mathbb{R}\) and
\(\nu_{\mathrm{def}}\) is supported at \(0\). This defect measure is
not the spectral measure of \(H_{\mathrm{HP}}\) and not the Li moment
measure; it records only the distance of zeros from the critical
line.
\end{remark}

\subsection{Li-Stieltjes Moment Formulation}
\label{v4-arith:vb:li-stieltjes}

\begin{proposition}[the defect measure is not a Li--Stieltjes moment model]
\label{v4-arith:vb:prop:li-stieltjes}
\ClaimStatusProvedHere. The off-line defect measure
\(\nu_{\mathrm{def}}\) of \Cref{v4-arith:vb:def:zero-measure} does
not give a Stieltjes moment representation of Li's coefficients or of
\(\mu_{n}:=\sum_{\rho}\rho^{n}\). A Li--Stieltjes formulation requires
the Hilbert--P\'olya zero-divisor operator \(\Theta_{\xi}\), the Li transform
\(\mathcal L_n\), and the positive Weil--Connes trace of
\Cref{v4-arith:rh:hyp:HPAr}.
\end{proposition}

\begin{proof}
\(\nu_{\mathrm{def}}\) records \(|\mathrm{Im}\,\gamma_{\rho}|^{2}\),
which vanishes exactly on the critical line. Li's summand is instead
\(1-(1-1/\rho)^{n}\), a holomorphic function of the zero
\(\rho\). No change of variables from the defect alone recovers this
holomorphic transform with multiplicity. The Hilbert--P\'olya trace identity
is \Cref{v4-arith:rh:prop:Li-is-Virasoro}, which applies the Li
transform to \(\Theta_{\xi}\).
\end{proof}

\subsection{Hamburger Moment Determinacy}
\label{v4-arith:vb:hamburger}

\begin{theorem}[Hamburger reformulation for the formal \(T_n\)]
\label{v4-arith:vb:thm:moment-equivalence}
\ClaimStatusProvedHere. The following equivalence is a formal
operator-theoretic template for a genuine positive moment problem; it
is not the live RH path unless the measure is replaced by the
Hilbert--P\'olya zero measure of \(H_{\mathrm{HP}}\) and the trace
is the positive Weil--Connes trace:
\begin{enumerate}
\item Formal positivity: \(T_{n}^{\mathrm{Vir}}\succeq0\) for all
\(n\geq1\) on a fixed self-adjoint spectral model.
\item The associated spectral measure is a positive measure supported
on \([1/2,\infty)\).
\item For the shifted moments
\(\nu_{k}:=\int_{[1/2,\infty)}(x-1/2)^{k}\,d\mu(x)\), all Hankel
minors \(H_{n}:=\det(\nu_{i+j})_{0\leq i,j\leq n}\) and shifted Hankel
minors \(H_{n}^{(1)}:=\det(\nu_{i+j+1})_{0\leq i,j\leq n}\) are
non-negative for every \(n\geq0\).
\end{enumerate}
This theorem is not an independent proof of RH and does not connect
\(T_{n}^{\mathrm{Vir}}\) to Li's coefficients.
\end{theorem}

\begin{proof}
\((1)\Leftrightarrow(2)\) is the spectral theorem for the chosen
self-adjoint model.  Since
\[
T_{2}^{\mathrm{Vir}}(x)=x^{2}-(x-1)^{2}=2x-1,
\]
positivity for all \(n\) forces \(x\geq1/2\).  Conversely, for
\(x=1/2+y\) with \(y\geq0\),
\((y+1/2)^{n}-(y-1/2)^{n}\) is non-negative.  This is a formal
statement about that model, not about zeta zeros.

\((2)\Leftrightarrow(3)\): Stieltjes moment-problem theorem
(Shohat--Tamarkin 1943, Ch.~III; Akhiezer 1965, \S2.1): a sequence
\(\{m_{k}\}_{k\geq 0}\) is the moment sequence of a positive measure on
\([0,\infty)\) iff all Hankel and shifted Hankel minors are
non-negative.  Apply it to the translated variable \(y=x-1/2\).

No implication to RH is asserted.
\end{proof}

\begin{remark}[moment information]
\label{v4-arith:vb:rem:moment-gain}
\Cref{v4-arith:vb:thm:moment-equivalence} records the formal
moment-determinacy apparatus that would become relevant after a true
Hilbert--P\'olya spectral measure is constructed. In particular:
\begin{itemize}
\item Testing a genuine positive spectral measure level-by-level
becomes testing Hankel
determinants \(H_{n}\) level-by-level, a finite numerical procedure at
each \(n\).
\item Known numerical bounds for Li coefficients do not test these
Hankel determinants; they test a different holomorphic transform of the
zero divisor.
\item Known structural estimates on Li coefficients and Voros
regularised sums do not supply lower bounds for the shifted Stieltjes
minors above.
\item The gap between known and needed becomes precisely specified:
\emph{unconditional} lower bounds on \(H_{n}^{(1)}\) for all \(n\).
\end{itemize}
The equivalent Hilbert--P\'olya condition is the construction of the
positive zero measure itself.
\end{remark}

\section{Local Kac Determinant at \(c=2\)}
\label{v4-arith:vb:kac}

\subsection{The Kac Formula}
\label{v4-arith:vb:kac-formula}

\begin{theorem}[Kac determinant formula]
\label{v4-arith:vb:thm:kac-formula}
\ClaimStatusProvedElsewhere (Kac 1978; Feigin--Fuchs 1990). For the
Virasoro Verma module \(M(c,h)\) of central charge \(c\) and highest
weight \(h\), the determinant of the contravariant form at level \(n\)
is
\begin{equation}
\label{v4-arith:vb:eq:kac-det}
\det_{n}^{c,h}\;=\;\prod_{\substack{r,s\geq 1\\rs\leq n}}(h-h_{r,s}(c))^{p(n-rs)}
\end{equation}
where \(p(k)\) is the partition function and
\begin{equation}
h_{r,s}(c)\;=\;\frac{((m+1)r-ms)^{2}-1}{4m(m+1)},\qquad c\;=\;1-\frac{6}{m(m+1)}.
\end{equation}
\end{theorem}

\subsection{The Arithmetic \(c=2\) Regime}
\label{v4-arith:vb:c=2-regime}

\begin{proposition}[Kac at \(c=2\)]
\label{v4-arith:vb:prop:kac-non-vanishing}
\ClaimStatusProvedHere. For the local rank-\(2\) free-boson Virasoro
module, \(c=2\). In the Kac parametrisation the parameter \(m\)
satisfies
\(1-\tfrac{6}{m(m+1)}=2\), hence \(m(m+1)=-6\), giving the formal
solution \(m=(-1\pm i\sqrt{23})/2\). The Kac null weights become
\begin{equation}
h_{r,s}(c=2)\;=\;\frac{((m+1)r-ms)^{2}-1}{4m(m+1)}\;=\;-\frac{((m+1)r-ms)^{2}-1}{24}.
\end{equation}
Thus the possible null weights are the countable Kac set
\(\{h_{r,s}(2):r,s\geq1\}\subset\mathbb{C}\). No rationalisation to
\(\{0,-1,-2,\ldots\}\) is implied.
\end{proposition}

\begin{proof}
Direct substitution \(c=2\) in the Kac formula gives the displayed
set. The free-boson realisation proves unitarity of the Fock module
itself; it does not identify the Kac zero locus with a real arithmetic
sequence.
\end{proof}

\begin{corollary}[no generic null-vector obstruction]
\label{v4-arith:vb:cor:no-generic-null}
\ClaimStatusProvedHere. For generic primary weight \(h\) outside the
countable Kac zero locus \(\{h_{r,s}(2):r,s\geq1\}\),
\(\det_{n}^{c=2,h}\neq 0\) at every level \(n\). This proves generic
non-degeneracy for the local Virasoro Verma module. It does not prove
Petersson positivity for the arithmetic primary object.
\end{corollary}

\subsection{Unitarity Gap at Special Weights}
\label{v4-arith:vb:unitarity-gap}

\begin{remark}[the gap between non-degeneracy and positivity]
\label{v4-arith:vb:rem:gap}
\Cref{v4-arith:vb:cor:no-generic-null} establishes non-degeneracy,
not positivity. Two Hermitian forms can be non-degenerate (no null
vectors) without being positive-definite (a form of signature
\((+,-,+,-,\ldots)\) is non-degenerate). The Friedan--Qiu--Shenker 1986
\(c<1\) analysis converts non-degeneracy to positivity via the
Kac determinant \emph{sign pattern}: when \(c<1\) is not in the
discrete unitary series, the sign of \(\det_{n}^{c,h}\) oscillates,
producing at least one negative eigenvalue and refuting unitarity.
The \(c=2\) case is outside the discrete \(c<1\) minimal unitary series,
but it lies in the standard \(c\geq1\) free-boson unitary regime. Hence
the FQS sign-oscillation argument for non-minimal \(c<1\) modules does
not apply.
\end{remark}

\begin{theorem}[\(c=2\) free-field unitarity]
\label{v4-arith:vb:thm:c=2-free-field-unitarity}
\ClaimStatusProvedElsewhere (Segal 1981; Frenkel--Ben-Zvi 2004).
The free-boson rank-\(2\) Fock module at central charge \(c=2\) is a
unitary representation of the Virasoro algebra, with Hermitian form
induced from the free-boson vacuum pairing
\(\langle\alpha_{-n_{1}}\cdots\alpha_{-n_{k}}|0\rangle,\alpha_{-m_{1}}\cdots\alpha_{-m_{l}}|0\rangle\rangle=\delta_{k,l}\sum_{\sigma\in S_{k}}\prod_{i}n_{i}\delta_{n_{i},m_{\sigma(i)}}\).
This is positive-definite by inspection.
\end{theorem}

\begin{corollary}[archimedean positivity of \(T_{n}^{\mathrm{Vir}}\)]
\label{v4-arith:vb:cor:archimedean-positivity}
\ClaimStatusProvedHere. On the \emph{archimedean} sector
\(\mathcal{H}^{(\infty)}\) of the primary arithmetic Heisenberg Fock
(the rank-\(2\) free-boson Fock of
\Cref{v4-arith:thm:P3-resolution}~Step~1), \(T_{n}^{\mathrm{Vir}}\succeq 0\) for every
\(n\geq 1\) on every non-vacuum number state.
\end{corollary}

\begin{proof}
The free-boson vacuum pairing of \Cref{v4-arith:vb:thm:c=2-free-field-unitarity}
is positive-definite; \(L_{0}^{\mathrm{Vir}}\) acts diagonally on
number states with eigenvalue \(m=\sum_{i}i n_{i}\geq 0\) (with \(m=0\)
precisely on the vacuum). On non-vacuum (\(m\geq 1\)) number states,
\(T_{n}^{\mathrm{Vir}}=(L_{0}^{\mathrm{Vir}})^{n}-(L_{0}^{\mathrm{Vir}}-\mathrm{id})^{n}\) acts by the scalar
\(m^{n}-(m-1)^{n}\geq 0\) (the function \(x\mapsto x^{n}\) is monotone on
\([0,\infty)\)). Positive-semidefiniteness follows.
\end{proof}

\begin{remark}[the missing sector]
\label{v4-arith:vb:rem:missing-sector}
\Cref{v4-arith:vb:cor:archimedean-positivity} establishes the
archimedean sector. The full primary arithmetic Heisenberg Fock is
\emph{not} just the archimedean sector: the Bost--Connes zero-mode
lattice \(L=\ell^{2}(\mathbb{N}^{\times})\)
(\Cref{v4-arith:def:zero-mode-lattice}) enters multiplicatively, and
the full primary space is
\(\mathcal{V}^{\mathrm{prim}}_{\mathrm{Heis},\mathrm{prim}}\cong L\otimes\mathcal{H}^{(\infty)}_{\mathrm{primary}}\).
Positivity of \(T_{n}^{\mathrm{Vir}}\) on the \(\mathcal{H}^{(\infty)}_{\mathrm{primary}}\)
factor is \Cref{v4-arith:vb:cor:archimedean-positivity}. Positivity
of a Virasoro polynomial on the \(L\)-factor is a statement about the
Bost--Connes energy spectrum \(\{\log m:m\geq1\}\). It is not a
statement about the zero divisor of \(\zeta\). The zero divisor enters
only through the Hilbert--P\'olya operator \(\Theta_{\xi}\) and the
Weil--Connes trace.
\end{remark}

\section{Fock-Combinatorial Identity}
\label{v4-arith:vb:fock-combinatorial}

\subsection{Explicit Action on Number States}
\label{v4-arith:vb:number-states}

\begin{lemma}[diagonal action on number states]
\label{v4-arith:vb:lem:diagonal-action}
\ClaimStatusProvedHere. On the Fock number state
\(|\mathbf{n}\rangle:=\prod_{i\geq 1}(\alpha_{-i}^{\infty})^{n_{i}}|0\rangle\big/\prod_{i}\sqrt{i^{n_{i}}n_{i}!}\)
(normalised with respect to the free-boson Fock inner product of
\Cref{v4-arith:vb:thm:c=2-free-field-unitarity}), \(L_{0}^{\mathrm{Vir}}\)
acts as
\begin{equation}
\label{v4-arith:vb:eq:L_0-action}
L_{0}^{\mathrm{Vir}}|\mathbf{n}\rangle\;=\;\Bigl(\sum_{i\geq 1}i\,n_{i}\Bigr)|\mathbf{n}\rangle\;=:\;m(\mathbf{n})|\mathbf{n}\rangle.
\end{equation}
\end{lemma}

\begin{proof}
\(L_{0}=\sum_{i\geq 1}\alpha_{-i}\alpha_{i}\) (Sugawara stress tensor
zero mode for the free boson; Frenkel--Ben-Zvi 2004~\S3.5). Applied
to \((\alpha_{-i})^{n_{i}}|0\rangle\): by the canonical commutator
\([\alpha_{i},\alpha_{-i}]=i\), \(\alpha_{i}(\alpha_{-i})^{n_{i}}|0\rangle=n_{i}\cdot i\cdot(\alpha_{-i})^{n_{i}-1}|0\rangle\),
hence \(\alpha_{-i}\alpha_{i}(\alpha_{-i})^{n_{i}}|0\rangle=i n_{i}(\alpha_{-i})^{n_{i}}|0\rangle\).
Summing over \(i\) gives the diagonal eigenvalue \(\sum_{i}i n_{i}\).
\end{proof}

\subsection{The Symmetric-Function Identity}
\label{v4-arith:vb:symmetric-identity}

\begin{theorem}[combinatorial positivity of \(T_{n}\) on Fock]
\label{v4-arith:vb:thm:fock-combinatorial-positivity}
\ClaimStatusProvedHere. On every non-vacuum number state
\(|\mathbf{n}\rangle\) with \(m(\mathbf{n})\geq 1\),
\begin{equation}
\label{v4-arith:vb:eq:fock-positivity}
T_{n}^{\mathrm{Vir}}|\mathbf{n}\rangle\;=\;[m(\mathbf{n})^{n}-(m(\mathbf{n})-1)^{n}]\,|\mathbf{n}\rangle,
\end{equation}
and \(m(\mathbf{n})^{n}-(m(\mathbf{n})-1)^{n}\geq 0\) for every
integer \(m(\mathbf{n})\geq 1\) and every \(n\geq 1\).
\end{theorem}

\begin{proof}
By \Cref{v4-arith:vb:lem:diagonal-action}, \(L_{0}^{\mathrm{Vir}}\) acts
by the scalar \(m=m(\mathbf{n})\). Hence \((L_{0}^{\mathrm{Vir}})^{n}\)
acts by \(m^{n}\) and \((L_{0}^{\mathrm{Vir}}-\mathrm{id})^{n}\) by
\((m-1)^{n}\). The difference is \(m^{n}-(m-1)^{n}\).

Non-negativity: the function \(x\mapsto x^{n}\) is monotone
non-decreasing on \([0,\infty)\) for \(n\geq 1\); hence for \(m\geq 1\) the
inequality \(m^{n}\geq (m-1)^{n}\) holds, with equality iff \(m=1\) and
\(n>0\) (in which case \(m^{n}-(m-1)^{n}=1\geq 0\)).
\end{proof}

\begin{remark}[combinatorial expansion]
\label{v4-arith:vb:rem:combinatorial-expansion}
Binomial expansion:
\begin{equation}
m^{n}-(m-1)^{n}\;=\;\sum_{k=0}^{n-1}\binom{n}{k}m^{k}(-1)^{n-1-k}\;=\;nm^{n-1}-\binom{n}{2}m^{n-2}+\cdots+(-1)^{n-1}.
\end{equation}
For \(m=1\): the identity \(1-0=1\); hence
\(T_{n}^{\mathrm{Vir}}|\mathbf{n}\rangle=|\mathbf{n}\rangle\)
on weight-\(1\) states for every \(n\geq 1\). This is the
\emph{singleton-primary}-state positivity.
\end{remark}

\begin{corollary}[Fock-level restricted positivity]
\label{v4-arith:vb:cor:fock-restricted-positivity}
\ClaimStatusProvedHere. On the subspace
\(\mathcal{F}^{m\geq 1}\subseteq\mathcal{H}^{(\infty)}\) spanned by
non-vacuum number states, \(T_{n}^{\mathrm{Vir}}\succeq 0\) for every \(n\geq 1\).
\end{corollary}

\begin{proof}
Direct from \Cref{v4-arith:vb:thm:fock-combinatorial-positivity}:
\(T_{n}^{\mathrm{Vir}}\) acts diagonally with non-negative eigenvalues; positivity
follows from the diagonal characterisation.
\end{proof}

\begin{remark}[zero-mode separation]
\label{v4-arith:vb:rem:fock-closed-not}
\Cref{v4-arith:vb:cor:fock-restricted-positivity} establishes
non-vacuum archimedean Fock positivity for \(T_{n}^{\mathrm{Vir}}\).
The primary arithmetic Heisenberg Fock also contains the
Bost--Connes zero-mode lattice
\(\ell^{2}(\mathbb{N}^{\times})\)
lattice where \(\mu_{n}\) acts via
\(\mu_{n}\delta_{m}=\delta_{nm}\); see
\Cref{v4-arith:thm:BC-acts-on-L}. On this lattice the Bost--Connes
Hamiltonian has Dirichlet-series spectrum
\(\{\log n\colon n\in\mathbb{N}^{\times}\}\) at KMS temperature
\(\beta\). This spectrum gives the Euler factor
\(\zeta(\beta)\); it is not the divisor of
\(\xi_{\mathbb R}\).
The RH-strength input is the separate trace bridge from this Euler
regularisation to \(\Theta_{\xi}\).
\end{remark}

\section{Reduction to the Irreducible Core Obstruction}
\label{v4-arith:vb:reduction-core}

\subsection{The Core Statement}
\label{v4-arith:vb:core-statement}

\begin{theorem}[conditional reduction to the zero-operator core]
\label{v4-arith:vb:thm:reduction}
\ClaimStatusProvedHere \emph{(conditional on P2, P3,
H-P\(^{\mathrm{Ar}}\), the primary tensor decomposition, and a positive
regularised trace)}. Under these hypotheses, the RH-strength
positivity statement reduces to the following zero-operator statement:
\begin{itemize}
\item[\textbf{(VB\(^{*}_{\mathrm{core}}\))}] \emph{Zero-operator spectral
positivity.} The operator
\(\Theta_{\xi}=1/2+iH_{\mathrm{HP}}\) satisfies
\(\det_{\infty}(s-\Theta_{\xi})=\xi_{\mathbb R}(s)\), and the
regularised Weil--Connes trace on the Li transforms
\(\mathcal L_n(\Theta_{\xi})\) is positive for every \(n\geq1\).
\end{itemize}
The Bost--Connes zero-mode lattice supplies the Euler/KMS
regularisation of the trace. It is not the spectrum of
\(\Theta_{\xi}\).
\end{theorem}

\begin{proof}
H-P\(^{\mathrm{Ar}}\) gives the determinant identity for
\(\Theta_{\xi}\). By
\Cref{v4-arith:rh:prop:Li-is-Virasoro}, the Li coefficient
\(\lambda_n\) is the regularised trace of
\(\mathcal L_n(\Theta_{\xi})\). Positivity of these traces for every
\(n\) is Li positivity, hence RH by Li's theorem. Conversely, RH is
the support condition for the spectral measure of \(\Theta_{\xi}\).
The Bost--Connes spectrum \(\{\log m\}\) enters only through the
regularisation giving the Euler product; no step identifies it with
the determinant divisor of \(\Theta_{\xi}\).
\end{proof}

\subsection{The Named Obstruction}
\label{v4-arith:vb:named-obstruction}

\begin{corollary}[the irreducible core of VB\(^{*}\)]
\label{v4-arith:vb:cor:irreducible-core}
\ClaimStatusProvedHere \emph{(conditional on the hypotheses of
\Cref{v4-arith:vb:thm:reduction})}. The irreducible core is the
positive zero-divisor measure problem for \(\Theta_{\xi}\). Replacing
\(\Theta_{\xi}\) by the Bost--Connes Hamiltonian \(N_{\mathrm{BC}}\)
replaces the zero divisor by the Euler eigenvalues and therefore
destroys the Hilbert--P\'olya statement.
\end{corollary}

\begin{remark}[core equivalence]
\label{v4-arith:vb:rem:core-meaning}
(VB\(^{*}_{\mathrm{core}}\)) is not a \emph{weaker} statement than
VB\(^{*}\) under the Hilbert--P\'olya and positive-trace hypotheses; it
is the zero-operator form of the same condition.
\begin{enumerate}
\item The zero divisor is carried by \(\Theta_{\xi}\).
\item The Euler product is carried by the Bost--Connes KMS trace.
\item The equality between the two is the Hilbert--P\'olya trace
bridge.
\item The Bost--Connes phase transition at \(\beta=1\) records the
pole of \(\zeta\); it is not a statement about phase transitions on
the critical line.
\end{enumerate}
\end{remark}

\section{Comparison to Classical RH Equivalents}
\label{v4-arith:vb:comparison}

\begin{table}[h]
\centering
\small
\begin{tabular}{l|l|l}
Formulation & Input data & Positivity shape \\
\hline
Riemann 1859 & \(\xi(s)\) entire of order \(1\) & Zero placement on line\\
Li 1997 & \(\lambda_{n}\), \(n\geq 1\) & Scalar sequence \(\geq 0\)\\
Weil 1952 & Test functions \(f\in\mathcal{S}\) & Inner product positivity\\
B\'aez-Duarte 2003 & Finite linear combinations of \(\{\rho^{s}\}\) & Density in \(L^{2}\)\\
Lagarias 1999 & de Branges datum over \(\Xi\) & Hermite--Biehler divisor control\\
Bombieri--Lagarias 1999 & \(\{\lambda_{n}\}\) with bounds & Explicit lower bounds\\
Connes 1999 & \(\mathcal{S}_{\mathbb{A}}/\mathbb{Q}^{*}\) & Trace-formula positivity\\
Vol~IV Virasoro & \(\{T_{n}^{\mathrm{Vir}}\}_{n\geq 1}\) in primary Fock & Fock polynomial positivity\\
\textbf{Vol~IV HP--Li core} & \(\boldsymbol{\Theta_{\xi}}\) with Weil--Connes trace & \textbf{zero-measure positivity}\\
\end{tabular}
\caption{RH-equivalent positivity formulations across the literature.}
\label{v4-arith:vb:tbl:comparison}
\end{table}

\begin{remark}[placement of VB\(^{*}\) against its neighbours]
\label{v4-arith:vb:rem:placement}
The formal Virasoro polynomial \(T_{n}^{\mathrm{Vir}}\) does not
imply Li's inequality. The RH-strength statement is HP--Li:
positivity of the regularised Weil--Connes trace of
\(\mathcal L_n(\Theta_{\xi})\) for every \(n\). Weil 1952 and the
de Branges formulations place the same zero-placement barrier in
classical languages. The Vol~IV core is the construction and
positivity of the \(\Theta_{\xi}\)-measure, with the Bost--Connes
zero-mode serving as Euler regularisation rather than as the zero
spectrum.
\end{remark}

\begin{proposition}[Weil comparison under P3]
\label{v4-arith:vb:prop:weil-equivalence}
\ClaimStatusProvedHere \emph{(conditional on Koszul--Petersson
compatibility \Cref{v4-arith:cl:cj:koszul-Petersson-compat} and the
H-P\(^{\mathrm{Ar}}\) trace condition)}. Weil 1952 explicit-formula
positivity and HP--Li positivity express the same RH-equivalent
barrier after identifying the Weil distribution with the regularised
trace of the zero-divisor operator \(\Theta_{\xi}\). The Petersson form enters
through the Koszul--Petersson comparison of the trace functional; it
does not turn \(T_{n}^{\mathrm{Vir}}\) into a Li transform.
\end{proposition}

\begin{proof}
Weil 1952 gives \(\lambda_{n}=\langle f_{n},f_{n}\rangle^{\mathrm{Weil}}\)
for \(f_{n}\) an explicit test function on the adele class space.
H-P\(^{\mathrm{Ar}}\) gives
\(\lambda_{n}=\mathrm{Tr}_{\mathrm{reg}}(\mathcal L_n(\Theta_{\xi}))\).
Koszul--Petersson compatibility identifies the trace pairing with the
Petersson pairing on the arithmetic primary sector whenever the
admissible Weil test lies in that sector. Thus Weil positivity and
HP--Li positivity are the same positivity problem after the trace
bridge is supplied. No equality with
\(\mathrm{Tr}_{\mathrm{reg}}(T_{n}^{\mathrm{Vir}})\) is used.
\end{proof}

\section{Zero-Divisor Condition}
\label{v4-arith:vb:residual-obstruction}

\subsection{The Precise Condition}
\label{v4-arith:vb:precise-residual}

The RH-strength positivity statement is equivalent to the following
condition:

\begin{definition}[the Vol~IV zero-divisor condition]
\label{v4-arith:vb:def:residual}
\ClaimStatusDefinitional. The \emph{Vol~IV zero-divisor condition} is
the following statement:
\begin{equation}
\label{v4-arith:vb:eq:residual}
\tag{\(\mathcal{R}_{\mathrm{IV}}\)}
\parbox{0.82\linewidth}{%
There is a normal zero-divisor operator
\(\Theta_{\xi}=1/2+iH_{\mathrm{HP}}\) whose determinant divisor is
\(\mathrm{div}\,\xi_{\mathbb R}\), and whose
Weil--Connes trace is supported on
\(1/2+i\mathbb R\). The Bost--Connes zero-mode gives the Euler/KMS
regularisation of the trace; it does not carry
\(\mathrm{div}\,\xi_{\mathbb R}\).%
}
\end{equation}
\end{definition}

\begin{theorem}[\(\mathcal{R}_{\mathrm{IV}}\) implies the Riemann Hypothesis]
\label{v4-arith:vb:thm:residual-is-RH}
\ClaimStatusProvedHere. \(\mathcal{R}_{\mathrm{IV}}\Rightarrow\) RH.
Conversely, RH gives only the zero-placement part of
\(\mathcal{R}_{\mathrm{IV}}\); it does not construct the normal
zero-divisor operator, the determinant divisor, the Weil--Connes
trace, or the Bost--Connes regularisation.
\end{theorem}

\begin{proof}
By definition, \(\mathcal{R}_{\mathrm{IV}}\) includes a
\(\Theta_{\xi}\)-determinant divisor and Weil--Connes trace
concentrated on the critical line. Since
\(\det_{\infty}(s-\Theta_{\xi})=\xi_{\mathbb R}(s)\), this is exactly
the zero-placement assertion for the non-trivial zeros of \(\zeta\),
hence RH by \Cref{v4-arith:rh:rem:equivalence-to-RH}. The converse
would require the extra operator-realization data named in the
statement; zero placement alone does not supply it.
\end{proof}

\subsection{Virasoro Boundary}
\label{v4-arith:vb:framework-change}

\begin{remark}[structural change]
\label{v4-arith:vb:rem:structural-change}
The arithmetic Virasoro reduction gives:
\begin{enumerate}
\item Closed \(n=1\) unconditionally
(\Cref{v4-arith:vb:thm:n=1}).
\item Reduced \(n=2\) to the formal Virasoro threshold
\(L_{0}^{\mathrm{Vir}}\succeq(1/2)\mathrm{id}\), not to zero placement
(\Cref{v4-arith:vb:cor:n=2-is-RH}).
\item Established Fock-level positivity for the archimedean sector
at every \(n\) (\Cref{v4-arith:vb:cor:fock-restricted-positivity}).
\item Reformulated the zero-measure problem as a Hamburger moment
determinacy problem for \(\Theta_{\xi}\)
(\Cref{v4-arith:vb:thm:moment-equivalence}).
\item Proved the Kac determinant at \(c=2\) is non-zero at generic
weights (\Cref{v4-arith:vb:prop:kac-non-vanishing}), establishing
absence of null-vector unitarity obstruction.
\item Placed HP--Li in correspondence with Weil positivity
(\Cref{v4-arith:vb:prop:weil-equivalence}) and the de Branges
Hermite-Biehler framework.
\item The RH-strength condition is a single zero-measure positivity
statement for \(\Theta_{\xi}\)
(VB\(^{*}_{\mathrm{core}}\),
\Cref{v4-arith:vb:thm:reduction}).
\end{enumerate}
VB\(^{*}\) itself remains open, as it is equivalent to RH.
\end{remark}

\begin{remark}[VB\(^{*}_{\mathrm{core}}\) boundary]
\label{v4-arith:vb:rem:proof-strategy}
The assertion (VB\(^{*}_{\mathrm{core}}\)) contains the following
independent analytic inputs:
\begin{enumerate}
\item A second KMS-positivity input beyond Bost--Connes: a constraint
on the zero-divisor measure forcing critical-line concentration.
\item An arithmetic-modular bootstrap input: a Kac-formula-style
analysis of the arithmetic Virasoro algebra on the Bost--Connes
\(C^{*}\)-algebra, extracting unitarity from null-vector absence.
\item An \(n\to\infty\) asymptotic input: a moment-generating-function
bootstrap strong enough to close the \(n\to\infty\) limit of
\(\lambda_{n}\) scaling.
\end{enumerate}
The first input belongs to the arithmetic site of Connes--Consani
arXiv:1405.4527. The second is the \(C^{*}\)-algebraic analogue of
Friedan--Qiu--Shenker. The third is an asymptotic moment problem.
\end{remark}

\section{Arithmetic Zhu Finiteness Test}
\label{v4-arith:vb:zhu-finiteness-test}

The tempting shortcut is Verlinde finite-dimensionality: prove that the
primary Bost--Connes core is an arithmetic Zhu-finite object, obtain a
finite-dimensional Hermitian modular functor, and read positivity from
ordinary linear algebra. The shortcut fails on the full zeta object.

\begin{definition}[Euler-faithful arithmetic Zhu quotient]
\label{v4-arith:vb:def:euler-faithful-zhu}
\ClaimStatusDefinitional. An \emph{Euler-faithful arithmetic Zhu
quotient} of the primary Bost--Connes core is a quotient \(Q\) of the
zero-mode lattice \(L=\ell^{2}(\mathbb{N}^{\times})\) satisfying:
\begin{enumerate}
\item the Bost--Connes operators descend:
\[
\mu_{n}\overline{\delta}_{m}=\overline{\delta}_{nm},\qquad
\mu_{n}^{*}\overline{\delta}_{m}=
\begin{cases}
\overline{\delta}_{m/n},& n\mid m,\\
0,& n\nmid m,
\end{cases}
\]
where \(\overline{\delta}_{m}\) is the image of
\(\delta_{m}\in L\);
\item the KMS Hamiltonian \(H=\log N\) descends with
\(H\overline{\delta}_{m}=(\log m)\overline{\delta}_{m}\) whenever
\(\overline{\delta}_{m}\neq0\);
\item the ordinary KMS character of the quotient is the full
Bost--Connes Euler product:
\[
\mathrm{Tr}_{Q}(e^{-\beta H})=\sum_{m\geq1}m^{-\beta}=\zeta(\beta)
\qquad(\mathrm{Re}\,\beta>1).
\]
\end{enumerate}
It is \emph{Zhu finite} if \(\dim_{\mathbb C}Q<\infty\).
\end{definition}

\begin{theorem}[no Euler-faithful Zhu finiteness]
\label{v4-arith:vb:thm:no-euler-faithful-zhu-finite}
\ClaimStatusProvedHere. No Euler-faithful arithmetic Zhu quotient of
the primary Bost--Connes core is Zhu finite.
\end{theorem}

\begin{proof}
If \(Q\) were finite-dimensional, the trace
\(\mathrm{Tr}_{Q}(e^{-\beta H})\) would be a finite sum
\(\sum_{j=1}^{r}a_{j}e^{-\beta E_{j}}\) after diagonalising the
finite-dimensional self-adjoint operator \(H\). In particular it extends to an
entire function of \(\beta\). Euler-faithfulness requires this finite
sum to equal \(\zeta(\beta)\) on the half-plane
\(\mathrm{Re}\,\beta>1\). By the identity theorem the meromorphic
continuation of \(\zeta\) would be entire, contradicting its simple
pole at \(\beta=1\). Hence \(Q\) is infinite-dimensional.

Equivalently, the faithful Bost--Connes action of
\Cref{v4-arith:thm:BC-acts-on-L} sends
\(\delta_{1}\) to \(\delta_{m}\) by \(\mu_{m}\delta_{1}=\delta_{m}\).
Keeping the full Euler product keeps one independent KMS eigenline at
every \(m\in\mathbb{N}^{\times}\). A finite Zhu quotient must kill or
identify infinitely many of these lines, and then it no longer has
character \(\zeta(\beta)\).
\end{proof}

\begin{corollary}[Verlinde finite-dimensionality cannot close full RH]
\label{v4-arith:vb:cor:verlinde-shortcut-fails}
\ClaimStatusProvedHere. Arithmetic Zhu/Verlinde finite-dimensionality
cannot prove VB\(^{*}_{\mathrm{core}}\) for the full
\(\mathcal{V}^{\mathrm{prim}}\) while preserving the completed zeta
partition function. The only Zhu-finite quotients compatible with the
Bost--Connes action are finite-prime or finite-conductor truncations;
their characters are finite Euler products or finite Dirichlet
polynomials, not \(\xi(s)\).
\end{corollary}

\begin{proof}
By \Cref{v4-arith:vb:thm:no-euler-faithful-zhu-finite}, a
finite-dimensional quotient cannot be Euler-faithful. A quotient
supported on a finite set \(S\) of primes has character
\(\prod_{p\in S}(1-p^{-\beta})^{-1}\), possibly with an archimedean
Tate factor. Such a finite Euler product is a finite-level arithmetic
modular-functor object; it is not the zeta object whose zero-placement
is RH. Therefore finite-dimensional Hermitian linear algebra can test
finite windows, but it cannot supply the missing positivity for the
full \(\Theta_{\xi}\)-spectral measure.
\end{proof}

\begin{remark}[central charge does not rescue Zhu finiteness]
\label{v4-arith:vb:rem:c-not-zhu-finite}
The residue value \(r_{\zeta}=2\) of
\Cref{v4-arith:thm:regularised-c} is a determinant-rank scalar, not a
finite rank and not a trace. Chapter~\ref{v4-arith:polyakov} proves
that the full infinite tensor has no global Virasoro action. Thus
\(r_{\zeta}=2\) cannot be used as a \(C_{2}\)-cofiniteness or
Zhu-finiteness input.
\end{remark}

\begin{remark}[surviving direction]
\label{v4-arith:vb:rem:zhu-surviving-direction}
The finite-dimensional construction survives only as a finite-window
technology: truncate to finitely many primes or to bounded conductor,
prove the corresponding arithmetic modular functor is finite, compute
its Hermitian form, and compare its limiting trace to the
Bost--Connes KMS trace. Any limiting theorem strong enough to recover
\(\mathrm{VB}^{*}_{\mathrm{core}}\) must add a new positivity or compactness
input beyond Zhu finite-dimensionality at finite level.
\end{remark}

\section{Zero-Measure Conjecture}
\label{v4-arith:vb:open-problem}

\begin{conjecture}[VB\(^{*}_{\mathrm{core}}\): the Vol~IV form of RH]
\label{v4-arith:vb:conj:VB-core}
\ClaimStatusConjectured. The zero-divisor operator
\(\Theta_{\xi}=1/2+iH_{\mathrm{HP}}\), with
\(\mathrm{div}\,\det_{\infty}(s-\Theta_{\xi})=\mathrm{div}\,\xi_{\mathbb R}\),
has determinant divisor and Weil--Connes trace supported on
\(1/2+i\mathbb R\).
Equivalently, all
non-trivial zeros of \(\zeta\) lie on the critical line.
\end{conjecture}

\begin{remark}[the conjecture is RH]
\label{v4-arith:vb:rem:conjecture-is-RH}
\Cref{v4-arith:vb:conj:VB-core} is, verbatim, the Riemann Hypothesis.
The arithmetic chiral-algebra formulation asks for a positive,
determinate Weil--Connes trace and determinant divisor on the zero-divisor operator
\(\Theta_{\xi}\). The Bost--Connes algebra supplies the Euler
regularisation and KMS weight, in the sense of Connes' adelic trace
formula; it does not supply the determinant divisor by itself.
\end{remark}

\begin{remark}[size of \( \mathrm{VB}^{*}_{\mathrm{core}} \)]
\label{v4-arith:vb:rem:substantial}
Any unconditional proof of \Cref{v4-arith:vb:conj:VB-core} would be
an unconditional proof of RH, a Millennium Prize problem. No
finite estimate is meaningful. Partial reductions have the following
mathematical sizes:
\begin{description}
\item[Bombieri-Lagarias-style sharpening at specific \(n\).]
substantial local theorem.
\item[Kac-determinant analysis at \(c=2\) to low level.]
finite-level determinant computation.
\item[Moment-problem Hankel analysis.] unbounded Hamburger
determinacy and lower bounds on all shifted Hankel minors.
\item[Arithmetic-site incorporation.] Connes--Consani arithmetic-site
positivity; no unconditional proof of RH follows.
\end{description}
\end{remark}

\section{Operator Boundary}
\label{v4-arith:vb:summary}

\begin{enumerate}
\item The formal Virasoro polynomial sequence
\(\{T_{n}^{\mathrm{Vir}}\}\) gives Fock positivity statements, not
Li transforms
(\Cref{v4-arith:vb:thm:three-equivalences}).
\item The RH-strength path uses
\(\mathcal L_n(\Theta_{\xi})=1-(1-\Theta_{\xi}^{-1})^n\), the
positive Weil--Connes trace, and Li's criterion
(\Cref{v4-arith:vb:thm:moment-equivalence},
\Cref{v4-arith:vb:prop:weil-equivalence}).
\item \(n=1\) is unconditional:
\(T_{1}^{\mathrm{Vir}}=\mathrm{id}\)
(\Cref{v4-arith:vb:thm:n=1}).
\item \(n=2\) reduces to
\(L_{0}^{\mathrm{Vir}}\succeq(1/2)\mathrm{id}\), a formal Fock
threshold rather than a zero-placement theorem
(\Cref{v4-arith:vb:thm:n=2},
\Cref{v4-arith:vb:cor:n=2-is-RH}).
\item Finite Li data at \(n\geq 3\) are scalar trace data for
\(\Theta_{\xi}\); they do not imply Virasoro operator positivity or
RH.
\item Fock-combinatorial positivity holds in the archimedean sector:
\(T_{n}^{\mathrm{Vir}}\succeq 0\) on non-vacuum number states for every \(n\)
(\Cref{v4-arith:vb:thm:fock-combinatorial-positivity}).
\item The local Kac determinant at \(c=2\) does not vanish at
generic weights, establishing absence of null-vector obstruction
(\Cref{v4-arith:vb:prop:kac-non-vanishing},
\Cref{v4-arith:vb:cor:no-generic-null}).
\item The RH-strength core is a single statement about the
\(\Theta_{\xi}\)-spectral measure: (VB\(^{*}_{\mathrm{core}}\),
\Cref{v4-arith:vb:thm:reduction}).
\item Arithmetic Zhu/Verlinde finite-dimensionality holds at finite windows
only: any Euler-faithful quotient preserving the
Bost--Connes character \(\zeta(\beta)\) is infinite-dimensional
(\Cref{v4-arith:vb:thm:no-euler-faithful-zhu-finite},
\Cref{v4-arith:vb:cor:verlinde-shortcut-fails}).
\item (VB\(^{*}_{\mathrm{core}}\)) is the Vol~IV form of RH,
equivalent to
\(\mathrm{Re}(\rho)=1/2\) for all non-trivial zeros
(\Cref{v4-arith:vb:thm:residual-is-RH}).
\item The condition \(\mathcal{R}_{\mathrm{IV}}\) is the existence of
the zero-divisor operator with positive, determinate Weil--Connes trace; it
is equivalent to RH.
\end{enumerate}

The unconditional Virasoro content is the identity at \(n=1\), the
formal threshold at \(n=2\), and non-vacuum Fock positivity. The
Riemann Hypothesis enters only through the zero-divisor operator
\(\Theta_{\xi}\), the Li transforms, and the positive Weil--Connes
trace. The equivalent condition is stated in two
operator-theoretic forms
(\Cref{v4-arith:vb:conj:VB-core},
\Cref{v4-arith:vb:def:residual}) and is, in full, the Riemann
Hypothesis.
